\chapter{The interface}

\section{The Now playing screen}
\capture{now-playing-en}{Now playing screen, information folded}

At the top of the screen, the toolbar holds the fullscreen switch, the live settings, the choice between the
oscilloscope and the filter response (for the SID), the dashboard toggle and the main menu. The heart of the screen
is the visualisation: one oscilloscope per voice, then the chip dashboard, described in
chapter~\ref{chap:dashboards}.

Below the visualisation, the information block gives the track's title, author, chip, format and original driver.
We fold it by default to leave the height to the visualisation; the arrow at the bottom left unfolds it. Then comes
the position in the track, then the row of controls: information, shuffle, previous, play and pause, next, repeat,
the $\equiv$+ button and the track menu.

The $\equiv$+ button is for filing the current track: you can add it to your favourites or to any playlist,
including a new one. The track menu, on the right, lets you share it, make it a ringtone or show all its properties.

\capture{now-playing-info-en}{Now playing screen, information unfolded}

\section{The Playlists screen}
\capture{playlist-en}{Playlists screen}
At the top, search filters the playlist shown; the buttons next to it give statistics, export, clear and sort. Just
below, the selector shows the current playlist and its number of tracks; tap it to switch playlists or manage them
(chapter~\ref{chap:listes}). Each row shows the chip family, title, author and duration. The handle on the right
moves a track.

\section{The Browser screen}
\capturedouble{browser-sources-en}{browser-folder-en}{Sources}{folder contents}
The path is shown at the top; tap one of its parts to jump straight back to that level. A long press on a file or
folder starts multiple selection.

\section{Away from the Now playing screen}
When you are in the playlists or the browser, a mini-player above the bottom bar reminds you of the current track
and lets you skip or pause without changing screens.

\section{Landscape, tablet, fullscreen}
\capture[0.9]{landscape-now-playing-en}{Now playing screen in landscape}
In landscape, the visualisation takes the width and the dashboard moves beside it. On a tablet, all three screens
make use of the available space.

\capture[0.9]{tablet-playlist-en}{On a tablet}

Fullscreen shows the oscilloscope alone; in portrait, the dashboard stays below it. We avoid the camera cutout at the
top and the rounded corners at the bottom in portrait, and the side cutout in landscape. A short tap leaves
fullscreen, a long press opens the live settings.
